\documentclass[11pt]{book}

\usepackage{amsmath, amssymb, amsthm}
\usepackage{geometry}
\usepackage{hyperref}
\usepackage{enumitem}

\geometry{margin=1in}

% theorem styles
\theoremstyle{definition}
\newtheorem{definition}{Definition}[chapter]
\newtheorem{example}[definition]{Example}
\newtheorem{remark}[definition]{Remark}
\newtheorem{insight}[definition]{Insight}

\theoremstyle{plain}
\newtheorem{theorem}[definition]{Theorem}
\newtheorem{proposition}[definition]{Proposition}
\newtheorem{corollary}[definition]{Corollary}

\newcommand{\R}{\mathbb{R}}
\newcommand{\C}{\mathbb{C}}
\newcommand{\Q}{\mathbb{Q}}
\newcommand{\Z}{\mathbb{Z}}
\newcommand{\N}{\mathbb{N}}

\title{Mathematical Analysis}
\author{Soumyajit De}
\date{May 21, 2021}

\begin{document}

\maketitle
\tableofcontents

% ------------------------------------------------
\chapter{Sets and Functions}

\section{Finite, Countable, and Uncountable Sets}

\subsection{Preliminary Definitions Related to Functions and Cardinality}

\begin{definition}[Function]
For two non-empty sets $A$ and $B$, if each element $x$ of $A$ is associated with an element of $B$,
then this association is called a function from $A$ to $B$, written as $f : A \to B$.
\end{definition}

\begin{insight}
Imagine packing when relocating: items go into boxes. A function specifies which item goes into which box.
Multiple items may go into the same box, but one item cannot go into two boxes.
\end{insight}

\begin{definition}[Domain]
The set $A$ is called the domain of $f$.
\end{definition}

\begin{definition}[Values]
The element in $B$ associated with $x \in A$ is called the value of $x$, written as $f(x)$.
\end{definition}

\begin{definition}[Co-domain]
The set $B$ is called the co-domain of $f$.
\end{definition}

\begin{definition}[Range]
The set $f(A) := \{ f(x) \mid x \in A \}$ is called the range of $f$.
\end{definition}

\begin{remark}
$f(A) \subseteq B$.
\end{remark}

\begin{definition}[Onto Mapping / Surjection]
If $f(A) = B$, then $f$ is called a mapping of $A$ onto $B$ or a surjection.
\end{definition}

\begin{definition}[Image]
For $E \subseteq A$, the set $f(E) := \{ f(x) \mid x \in E \}$ is called the image of $E$.
\end{definition}

\begin{definition}[Pre-image]
For $E \subseteq B$, the set
\[
f^{-1}(E) := \{ x \in A \mid f(x) \in E \}
\]
is called the pre-image of $E$.
\end{definition}

\begin{definition}[One-to-one Mapping / Injection]
$f$ is injective if for every $y \in f(A)$, the set $f^{-1}(\{y\})$ consists of exactly one element.
\end{definition}

\begin{definition}[Bijection]
If there exists a one-to-one mapping of $A$ onto $B$, then $A$ and $B$ are said to be in one-to-one correspondence,
written as $A \sim B$.
\end{definition}

\begin{remark}
If $A \sim B$, then $A$ and $B$ have the same cardinal number.
\end{remark}

\subsection{Finite, Countable, and Uncountable Sets}

Let $J_n := \{1, \dots, n\}$ and $J := \{1,2,3,\dots\}$.

\begin{definition}[Finite Set]
A set $A$ is finite if $A = \emptyset$ or $A \sim J_n$ for some $n$.
\end{definition}

\begin{definition}[Infinite Set]
A set that is not finite is infinite.
\end{definition}

\begin{definition}[Countable Set]
If $A \sim J$, then $A$ is countable.
\end{definition}

\begin{definition}[Uncountable Set]
A set that is neither finite nor countable is uncountable.
\end{definition}

\begin{definition}[Sequence]
A function $f : J \to X$ is called a sequence, usually written as $\{x_n\}$.
\end{definition}

\subsection{Properties}

\begin{theorem}
Every infinite subset of a countable set is countable.
\end{theorem}

\begin{theorem}
A countable union of countable sets is countable.
\end{theorem}

\begin{theorem}
The set of all sequences whose elements are taken from a finite set is uncountable.
\end{theorem}

% ------------------------------------------------
\chapter{The Real and Complex Number Systems}

\section{Groups}

\begin{definition}[Group]
A group $(G,+,0)$ is a set $G$ with a binary operation $+$ satisfying closure, associativity,
existence of identity, and existence of inverses.
\end{definition}

\begin{definition}[Abelian Group]
A group is abelian if $x+y=y+x$ for all $x,y \in G$.
\end{definition}

\section{Fields}

\begin{definition}[Field]
A field $(F,+,0,\cdot,1)$ is a set with two operations satisfying the group axioms and distributivity.
\end{definition}

\section{Ordered Sets}

\begin{definition}[Ordered Set]
An ordered set is a set equipped with a total order relation.
\end{definition}

\begin{definition}[Least Upper Bound Property]
An ordered set has the least upper bound property if every nonempty bounded-above subset
has a supremum.
\end{definition}

\section{The Real Field}

\begin{definition}[Real Numbers]
There exists an ordered field $\R$ with the least upper bound property containing $\Q$ as a subfield.
\end{definition}

\section{The Complex Field}

\begin{definition}[Complex Numbers]
$\C$ consists of ordered pairs $(a,b)$ with addition and multiplication defined by
\[
(a,b)+(c,d)=(a+c,b+d), \quad (a,b)(c,d)=(ac-bd,ad+bc).
\]
\end{definition}

% ------------------------------------------------
\chapter{Basic Topology}

\section{Metric Spaces}

\begin{definition}[Metric Space]
A metric space is a set $X$ with a function $d:X\times X\to \R_{\ge 0}$ satisfying positivity,
symmetry, and the triangle inequality.
\end{definition}

\begin{definition}[Open Set]
A set $E$ is open if every point of $E$ is an interior point.
\end{definition}

\begin{definition}[Closed Set]
A set is closed if it contains all of its limit points.
\end{definition}

\section{Compact Sets}

\begin{definition}[Compactness]
A set is compact if every open cover admits a finite subcover.
\end{definition}

\begin{theorem}[Heine--Borel]
A subset of $\R^k$ is compact if and only if it is closed and bounded.
\end{theorem}

% ------------------------------------------------
\chapter{Numerical Sequences and Series}

\begin{definition}[Convergent Sequence]
A sequence $\{p_n\}$ converges to $p$ if for every $\varepsilon>0$ there exists $N$ such that
$n>N$ implies $d(p_n,p)<\varepsilon$.
\end{definition}

\begin{definition}[Cauchy Sequence]
A sequence is Cauchy if its terms eventually become arbitrarily close to each other.
\end{definition}

\begin{theorem}
Every Cauchy sequence in $\R^k$ converges.
\end{theorem}

\section{Upper and Lower Limits}

\begin{definition}[$s_n \to +\infty$]
Let $\{s_n\}$ be a sequence of real numbers such that for every $M \in \R$, there exists
$N \in \Z^+$ such that $n > N \Rightarrow s_n > M$.
We write $s_n \to +\infty$.
\end{definition}

\begin{definition}[$s_n \to -\infty$]
Let $\{s_n\}$ be a sequence of real numbers such that for every $M \in \R$, there exists
$N \in \Z^+$ such that $n > N \Rightarrow s_n \le M$.
We write $s_n \to -\infty$.
\end{definition}

\begin{remark}
Difference in notation between $\{s_n\} \to p$ and $s_n \to +\infty$.
\end{remark}

\begin{remark}
$\{s_n\} \to p$ is a convergent sequence whereas $s_n \to \pm\infty$ are divergent sequences.
\end{remark}

\begin{definition}[Upper and Lower Limits]
Let $\{s_n\}$ be a sequence of real numbers. Let $E$ be the set of numbers
$x \in \R \cup \{-\infty,+\infty\}$ such that $s_{n_k} \to x$ for some subsequence $\{s_{n_k}\}$.
Define
\[
s^* = \sup E, \quad s_* = \inf E.
\]
Then $s^*$ is called the upper limit and $s_*$ the lower limit of $\{s_n\}$.
We write
\[
\limsup_{n\to\infty} s_n = s^*, \qquad
\liminf_{n\to\infty} s_n = s_*.
\]
\end{definition}

\begin{theorem}
For a sequence $\{s_n\}$ of real numbers:
\begin{enumerate}
\item $s^*, s_* \in E$.
\item If $x > s^*$, then there exists $N$ such that $n > N \Rightarrow s_n < x$.
\item If $x < s_*$, then there exists $N$ such that $n > N \Rightarrow s_n > x$.
\end{enumerate}
\end{theorem}

\begin{theorem}
If $s_n \le t_n$ for $n > N$, then
\[
\liminf_{n\to\infty} s_n \le \liminf_{n\to\infty} t_n,
\qquad
\limsup_{n\to\infty} s_n \le \limsup_{n\to\infty} t_n.
\]
\end{theorem}

\section{Some Special Sequences}

The following limits are commonly used:
\begin{enumerate}
\item $\displaystyle \lim_{n\to\infty} \frac{1}{n^p} = 0$, where $p>0$.
\item $\displaystyle \lim_{n\to\infty} n^{1/p} = 1$, where $p>0$.
\item $\displaystyle \lim_{n\to\infty} n^{1/n} = 1$.
\item $\displaystyle \lim_{n\to\infty} (n!)^{1/n} = +\infty$.
\item $\displaystyle \lim_{n\to\infty} a^n = 0$, where $|a|<1$.
\item $\displaystyle \lim_{n\to\infty} a^{1/n} = 1$.
\item $\displaystyle \lim_{n\to\infty} \left(1+\frac{1}{n}\right)^n = e$.
\item $\displaystyle \lim_{n\to\infty} \left(1+\frac{a}{n}\right)^{bn} = e^{ab}$.
\item $\displaystyle \lim_{n\to\infty} \arctan n = \frac{\pi}{2}$.
\item $\displaystyle \lim_{n\to\infty} \frac{n^\alpha}{a^n} = 0$, where $a>1$.
\item $\displaystyle \lim_{n\to\infty} \frac{P(n)}{e^n} = 0$, where $P(n)$ is a polynomial.
\item $\displaystyle \lim_{n\to\infty} \frac{\log n}{n^a} = 0$, where $a>0$.
\item $\displaystyle \lim_{n\to\infty} \frac{a^n}{n!} = 0$.
\end{enumerate}

\section{Series}

\section{Series of Nonnegative Terms}

\section{The Number \texorpdfstring{$e$}{e}}

\section{The Root and Ratio Tests}

\section{Power Series}

\section{Summation by Parts}

\section{Absolute Convergence}

\section{Addition and Multiplication of Series}

\section{Rearrangements}

% ------------------------------------------------
\chapter{Continuity}

\begin{definition}[Continuity]
A function $f:X\to Y$ is continuous at $p$ if for every $\varepsilon>0$ there exists $\delta>0$
such that $d_X(x,p)<\delta$ implies $d_Y(f(x),f(p))<\varepsilon$.
\end{definition}

\begin{theorem}[Intermediate Value Theorem]
A continuous function on an interval attains every intermediate value.
\end{theorem}

\begin{theorem}
A continuous function on a compact set is uniformly continuous.
\end{theorem}

\section{Limit of a Function}

\begin{definition}[Limit of a Function]
Let $E \subset X$, $f:E\to Y$, and $p \in E^\uparrow$.
If there exists $q \in Y$ such that for every $\varepsilon>0$
there exists $\delta>0$ such that
\[
0 < d_X(x,p) < \delta \Rightarrow d_Y(f(x),q) < \varepsilon,
\]
then $\lim_{x\to p} f(x)=q$.
\end{definition}

\begin{remark}
$p$ need not belong to the domain of $f$.
\end{remark}

\begin{remark}
Even if $f(p)$ exists, it is possible that $\lim_{x\to p} f(x)\ne f(p)$.
\end{remark}

\begin{remark}
If $E^\uparrow=\varnothing$, then $f$ has no limits.
\end{remark}

\begin{theorem}[Sequential Criterion]
$\lim_{x\to p} f(x)=q$ iff for every sequence $\{p_n\}\subset E$ with
$p_n\ne p$ and $p_n\to p$, we have $f(p_n)\to q$.
\end{theorem}

\section{Continuity}

\begin{definition}[Continuity]
$f:E\to Y$ is continuous at $p\in E$ if for every $\varepsilon>0$
there exists $\delta>0$ such that
\[
d_X(x,p)<\delta \Rightarrow d_Y(f(x),f(p))<\varepsilon.
\]
\end{definition}

\begin{theorem}
If $p$ is a limit point of $E$, then $f$ is continuous at $p$
iff $\lim_{x\to p} f(x)=f(p)$.
\end{theorem}

\section{Continuity and Compactness}

\begin{theorem}
If $f:X\to Y$ is continuous and $X$ is compact, then $f(X)$ is compact.
\end{theorem}

\begin{theorem}
If $f:X\to Y$ is a continuous bijection and $X$ is compact,
then $f^{-1}$ is continuous.
\end{theorem}

\begin{theorem}
If $f:X\to\R$ is continuous and $X$ is compact, then $f$ attains
its maximum and minimum.
\end{theorem}

\section{Uniform Continuity}

\begin{definition}[Uniform Continuity]
$f:X\to Y$ is uniformly continuous if for every $\varepsilon>0$
there exists $\delta>0$ such that
\[
d_X(p,q)<\delta \Rightarrow d_Y(f(p),f(q))<\varepsilon
\]
for all $p,q\in X$.
\end{definition}

\begin{theorem}
If $X$ is compact, then every continuous function $f:X\to Y$
is uniformly continuous.
\end{theorem}

\section{Lipschitz Continuity}

\begin{definition}[Lipschitz Continuity]
$f:X\to Y$ is Lipschitz continuous if there exists $K>0$ such that
\[
d_Y(f(p),f(q))\le K\,d_X(p,q)
\]
for all $p,q\in X$.
\end{definition}

\begin{theorem}
Every Lipschitz continuous function is uniformly continuous.
\end{theorem}

\section{Connectedness}

\begin{theorem}
If $f:X\to Y$ is continuous and $E\subset X$ is connected,
then $f(E)$ is connected.
\end{theorem}

\begin{theorem}[Intermediate Value Theorem]
If $f:[a,b]\to\R$ is continuous, then it attains every value between
$f(a)$ and $f(b)$.
\end{theorem}

\section{Monotonic Functions}

\begin{theorem}
If $f:(a,b)\to\R$ is monotonic, then the set of its discontinuities
is at most countable.
\end{theorem}


\end{document}
